\chapter{APIs, data modelling, and persistence}

\section{An interface is a promise over time}
An \term{API} is more than a URL or method signature. It defines inputs, outputs,
errors, identity, authorization, ordering, consistency, and compatibility. A
client can depend on behaviour that was never written down, so an API should
make its important promises deliberate.

For an HTTP endpoint, document at least:

\begin{itemize}
  \item the resource and method semantics;
  \item accepted representations and validation rules;
  \item success and error status codes with stable error shapes;
  \item authentication, authorization, rate limits, and sensitive fields;
  \item retry and idempotency behaviour;
  \item pagination, ordering, time zones, and version-compatibility rules.
\end{itemize}

HTTP is standardized as a stateless application-level protocol with shared
semantics across versions; the standard does not dictate your business resource
model \cite{ietf2022http}. This distinction matters. A `POST /orders` contract
is yours, while the semantics of status codes and method properties are governed
by the HTTP specification.

\section{Identifiers and time}
Identifiers should communicate less than they need to. An opaque identifier
reduces accidental coupling to storage order and makes enumeration harder, but
it does not replace authorization. If a client needs chronological ordering,
provide an explicit timestamp or cursor rather than asking it to decode an ID.

Time has at least three meanings: an instant in a global timeline, a local
calendar representation, and a duration measured by a monotonic clock. Store
instants in an unambiguous form, carry time-zone information when presenting
calendar values, and use a monotonic clock for elapsed-time measurements. Never
use wall-clock time to decide whether a timeout duration has elapsed.

\section{Relational modelling}
A relational schema is an executable set of invariants. In the example schema,
positive quantities are checked in the database even though application code
also validates them. Defence in depth is useful because imports, migrations,
administrative scripts, and future services can bypass one validation path.

The order total is stored as integer cents. Decimal money is a domain choice;
the important point is that binary floating-point approximations are not silently
used for a value whose legal meaning is exact. The database view recalculates the
sum of line items. In a production design, decide whether the stored total is a
snapshot, a cache, or an authoritative ledger value, and reconcile it explicitly.

\begin{lstlisting}[language=SQL,caption={A constraint is part of the data contract.}]
CREATE TABLE order_items (
    order_id TEXT NOT NULL REFERENCES orders(id),
    sku TEXT NOT NULL CHECK (length(trim(sku)) > 0),
    quantity INTEGER NOT NULL CHECK (quantity > 0),
    unit_price_cents INTEGER NOT NULL CHECK (unit_price_cents > 0),
    PRIMARY KEY (order_id, sku)
);
\end{lstlisting}

\section{Transactions and consistency}
A transaction groups changes so that the database presents an allowed state at
its commit boundary. The exact isolation guarantees vary by database engine and
configuration. Do not infer serializable behaviour from the word “transaction.”
Ask what anomalies are possible: dirty reads, lost updates, write skew, or
phantoms. Then choose a constraint, lock, isolation level, or retry strategy that
matches the invariant.

The outbox pattern is useful when a durable state change must cause a message.
Write the domain change and an outbox row in the same local transaction. A
separate publisher reads unclaimed outbox rows and sends them with at-least-once
delivery. Consumers must therefore be idempotent. The pattern does not create
exactly-once effects across arbitrary systems; it creates a recoverable handoff.

\section{Schema evolution}
A safe migration is compatible with the versions that may coexist during rollout.
The expand-and-contract sequence is a common strategy:

\begin{enumerate}
  \item add a nullable or additive representation;
  \item deploy code that can read both and writes the new form;
  \item backfill and validate existing data;
  \item switch readers and enforce the new invariant;
  \item remove the old representation only after rollback is no longer needed.
\end{enumerate}

This costs temporary complexity. The alternative, a single destructive change,
may be reasonable for a private database with a verified restore path and a
maintenance window. The decision depends on data value, downtime tolerance,
rollout topology, and rollback feasibility.

\begin{exercise}
Design an endpoint for retrieving a customer's orders. Specify pagination,
ordering, authorization, deletion behaviour, and what happens when a client
repeats a request after a timeout.
\end{exercise}

\section{Chapter summary}
APIs are long-lived behavioural contracts. Model identity, time, errors, retries,
and compatibility explicitly. Use database constraints as executable invariants,
understand the actual transaction guarantees, and evolve schemas in a way that
respects the versions present during deployment.

